\documentclass[10pt,conference,compsocconf]{IEEEtran}

\usepackage[hidelinks]{hyperref}
\usepackage{graphicx}
\usepackage{amsmath}
\usepackage{booktabs}

\begin{document}
\title{Predicting Coronary Heart Disease from the BRFSS Survey\\with Regularized Linear Models}

\author{
  Swapnil Bag, Sam Lee, Valentin Marbach\\
  \textit{CS-433 Machine Learning, EPFL, Switzerland}
}

\maketitle

\begin{abstract}
We predict whether a person has coronary heart disease (MICHD) from 321
answers to the 2015 Behavioral Risk Factor Surveillance System (BRFSS) survey,
using only linear models implemented from scratch with NumPy. The data is
heavily imbalanced (8.8\% positives), mostly missing (45\% of the entries) and
full of survey-specific codes. We build a preprocessing pipeline driven by an
exploratory analysis, evaluate every choice with stratified 5-fold
cross-validation and the F1 score, and compare the six required methods. A
logistic regression with a tuned step size and an F1-optimal decision
threshold reaches a cross-validated F1 of 0.431 and an F1 of 0.440 (accuracy
0.875) on the AIcrowd test set.
\end{abstract}

\section{Introduction}

Cardiovascular disease is the leading cause of death worldwide, and
self-reported health surveys are a cheap source of risk information. The task
is a binary classification problem: 328{,}135 training respondents with a
label \texttt{\_MICHD} in $\{-1,1\}$ and 109{,}379 test respondents. Two
properties of the data drive the whole approach. First, only 8.8\% of the
respondents are positive: predicting ``no disease'' for everyone gives 91.2\%
accuracy but is useless, so we optimize and report the F1 score (predicting
everyone positive gives $F_1 = 0.162$, our floor). Second, the features are raw
survey answers, with special codes, skipped questions and redundant derived
variables, so careful preprocessing is needed before any linear model.

We contribute an exploratory analysis that motivates each preprocessing
step, a leak-free evaluation protocol, an ablation study, and a comparison of
the six required methods with step sizes derived from the loss curvature.

\section{Data Analysis}
\label{sec:data}

\textbf{Missing values.} 44.8\% of all entries are missing; 99 features are
missing for more than 90\% of the respondents (optional modules such as vision
or asthma). Missingness is informative: some questions are only asked
depending on a previous answer. For instance, the diabetes module is only
asked to diabetics, and the positive rate is 22\% when it is filled versus 8\%
when it is missing.

\textbf{Special codes.} BRFSS encodes ``don't know'' as 7, 77, 777, ...,
``refused'' as 9, 99, ..., and ``none'' as 88 or 888 (e.g.\
\texttt{PHYSHLTH}~$=88$ means zero days of poor health)~\cite{codebook}. Left
as numbers, a linear model reads ``refused'' (9) as worse than ``poor
health'' (5). We detect these codes automatically: a candidate code is
flagged in a feature when it appears while the value just below it does not
(the code is isolated from the regular range).

\textbf{Useless and redundant features.} Six features are constant, and the
survey administration features (state, interview date, record numbers,
sampling weights) describe the interview rather than the person. Raw weight
and height mix imperial and metric units (metric values are stored as
$9000+$value), while the CDC provides clean computed versions (weight in kg,
height in m, BMI). We also found 26 pairs of near-duplicate features
($|\rho| > 0.95$), such as three encodings of age.

\textbf{Signal.} The known cardiovascular risk factors are clearly visible:
the positive rate grows from 0.5\% (age 18--24) to 22\% (80+) and from 2\%
(excellent general health) to 32\% (poor), and reaches 37\% after a stroke.
However, the strongest single correlation with the target is only about
0.25: the model has to combine many weak signals.

\section{Methods}

\subsection{Evaluation protocol}

All results use the same stratified 5-fold cross-validation (seed 1): each
fold keeps the 8.8\% positive rate, and every respondent is scored once by a
model that never saw it. The preprocessing is fitted on the four training
folds only, so no statistic of the held-out fold leaks into training. Because
of the imbalance, we do not threshold at 0.5: the decision threshold is the
one maximizing F1 on the out-of-fold scores (one sort and a cumulative sum
give the F1 of every threshold). We report the mean and standard deviation of
F1 over the folds.

\subsection{Preprocessing}

Following Section~\ref{sec:data}, the pipeline applies, in order:
(1)~drop the administration and raw mixed-unit features;
(2)~map special codes, ``none'' to 0 and ``don't know''/``refused'' to
missing;
(3)~add a 0/1 indicator for each feature that is sometimes missing (identical
indicators are merged);
(4)~drop constant features and features missing for more than 90\% of the
training rows;
(5)~drop the second feature of each pair with $|\rho| > 0.95$;
(6)~impute the remaining missing values with the training median;
(7)~one-hot encode integer features with 3 to 15 distinct values;
(8)~standardize every column and add a bias.
The 321 raw features become 619 columns: 125 numerical, 291 one-hot and 203
missing-value indicators.

\subsection{Models}

We use the six required methods: gradient descent (GD) and stochastic
gradient descent (SGD, one sample per step) on the mean squared error (MSE),
least squares, ridge regression, and logistic regression without and with
$L_2$ regularization. Labels are mapped to $\{0,1\}$; linear models output a
real score and logistic models a probability. Least squares solves the normal
equations with a pseudo-inverse, because one-hot columns of the same feature
can be linearly dependent.

\textbf{Step sizes.} GD on a smooth loss converges for $\gamma < 2/L$, with
$L$ the largest eigenvalue of the Hessian. For the MSE,
$L = \lambda_{\max}(X^\top X/N) \approx 42$ on our features, hence
$\gamma < 0.047$. For the logistic loss the Hessian is bounded by
$X^\top X/(4N)$, hence $\gamma < 0.19$. For single-sample SGD, a step is
stable only if $\gamma < 2/\|x_n\|^2$ for every row; most rows have
$\|x_n\|^2 \approx 500$ but the largest reaches about $330{,}000$, so any
$\gamma$ above $6\cdot 10^{-6}$ eventually diverges. We tune $\gamma$, the
number of iterations and $\lambda$ inside these ranges.

\section{Results}

\textbf{Preprocessing ablation.} Table~\ref{tab:ablation} removes one step at
a time (ridge, $\lambda = 10^{-6}$). The full pipeline improves F1 from 0.412
(mean imputation and scaling only) to 0.420. Removing any single step changes
F1 by less than the fold-to-fold standard deviation: with a linear model, the
gain comes from the combination of steps, one-hot encoding and missing
indicators being the most useful. Mapping ``don't know'' to missing slightly
lowers F1, which suggests that not knowing an answer is itself informative.

\begin{table}[htbp]
  \centering
  \caption{Preprocessing ablation (ridge, 5-fold F1).}
  \label{tab:ablation}
  \small\setlength{\tabcolsep}{4pt}
  \begin{tabular}{lcc}
    \toprule
    Variant & F1 & Prec. / Rec. \\
    \midrule
    All steps & $0.420 \pm 0.008$ & 0.338 / 0.555 \\
    \quad w/o admin drop & $0.420 \pm 0.008$ & 0.341 / 0.547 \\
    \quad w/o special codes & $0.422 \pm 0.005$ & 0.355 / 0.520 \\
    \quad w/o missing indicators & $0.418 \pm 0.006$ & 0.334 / 0.558 \\
    \quad w/o sparse drop & $0.423 \pm 0.007$ & 0.338 / 0.564 \\
    \quad w/o redundancy drop & $0.421 \pm 0.005$ & 0.351 / 0.526 \\
    \quad w/o one-hot & $0.416 \pm 0.007$ & 0.334 / 0.552 \\
    Impute + scale only & $0.412 \pm 0.004$ & 0.326 / 0.560 \\
    \bottomrule
  \end{tabular}
\end{table}

\textbf{Comparison of the methods.} Table~\ref{tab:models} reports each
method with its best setting. All methods share the full preprocessing and the
same folds. The squared-loss methods (GD, SGD, least squares, ridge) reach the
same solution and an F1 of about 0.420; GD converges within 1000 iterations
($\gamma = 0.04$) and SGD matches it with $\gamma = 5\cdot 10^{-6}$ and
3 million steps, whereas $\gamma = 10^{-5}$ diverges, as predicted above.
Ridge is insensitive to $\lambda$ up to $10^{-2}$ and degrades for
$\lambda = 1$. Logistic regression is the best method: with the initial step
size $\gamma = 0.01$ it had not converged (F1 0.422); $\gamma = 0.15$,
close to the theoretical bound, converges in about 2000 iterations and reaches
$0.431 \pm 0.005$, a gain larger than the fold-to-fold variation.
Regularization does not help ($\lambda \geq 10^{-2}$ lowers F1): with
328{,}000 rows for 620 parameters, the model hardly overfits.

\begin{table}[htbp]
  \centering
  \caption{Best setting of each method (5-fold F1).}
  \label{tab:models}
  \small\setlength{\tabcolsep}{4pt}
  \begin{tabular}{llc}
    \toprule
    Method & Setting & F1 \\
    \midrule
    Always positive & -- & $0.162 \pm 0.000$ \\
    MSE GD & $\gamma = 0.04$, 1000 it. & $0.420 \pm 0.007$ \\
    MSE SGD & $\gamma = 5\cdot10^{-6}$, 3M it. & $0.419 \pm 0.007$ \\
    Least squares & -- & $0.420 \pm 0.008$ \\
    Ridge & $\lambda = 10^{-4}$ & $0.420 \pm 0.008$ \\
    Logistic & $\gamma = 0.15$, 2000 it. & $\mathbf{0.431 \pm 0.005}$ \\
    Reg. logistic & $\lambda{=}10^{-3}$, $\gamma{=}0.15$, 2000 & $0.431 \pm 0.006$ \\
    \bottomrule
  \end{tabular}
\end{table}

\textbf{Final model.} We retrain the logistic regression ($\gamma = 0.15$,
2000 iterations) on all the training data and predict positive when the
probability exceeds the cross-validated threshold 0.20. The low threshold
compensates for the class imbalance: it trades accuracy for recall
(precision 0.355, recall 0.549 in cross-validation). On the AIcrowd test set
the submission obtains an \textbf{F1 of 0.440} and an accuracy of 0.875,
consistent with the cross-validated estimate (0.431), which shows that the
protocol did not overfit the threshold or the hyperparameters.

\section{Discussion}

The main limitation is the linearity of the models: all methods plateau
between 0.42 and 0.43, and neither regularization nor any single
preprocessing step moves this ceiling. Interactions between risk factors
(e.g.\ age with general health), polynomial features of the ordinal answers,
and keeping ``don't know'' as a separate category are natural next steps.
The threshold and hyperparameters are selected on the out-of-fold scores we
report, but the test score shows that this bias is small.

\section{Summary}

Data-driven preprocessing and a leak-free, imbalance-aware evaluation matter
as much as the model. A logistic regression with a curvature-based step size
and an F1-optimal threshold reaches F1 0.440 on the test set; \texttt{run.py}
in our repository regenerates this submission exactly.

\begin{thebibliography}{1}
\bibitem{codebook}
Centers for Disease Control and Prevention, ``Behavioral Risk Factor
Surveillance System: 2015 codebook report,''
\url{https://www.cdc.gov/brfss/annual_data/2015/pdf/codebook15_llcp.pdf}.
\end{thebibliography}

\end{document}
